\documentclass[11pt,letterpaper]{article}

\usepackage[T1]{fontenc}
\usepackage[utf8]{inputenc}
\usepackage{lmodern}
\usepackage[margin=0.92in,headheight=15pt]{geometry}
\usepackage{amsmath,amssymb,amsthm,mathtools,mathrsfs}
\usepackage{microtype}
\usepackage{booktabs,longtable,tabularx,array}
\usepackage{enumitem}
\usepackage[dvipsnames]{xcolor}
\usepackage{fancyhdr}
\usepackage{needspace}
\usepackage{xurl}
\usepackage[colorlinks=true,linkcolor=MidnightBlue,citecolor=MidnightBlue,
  urlcolor=MidnightBlue,breaklinks=true]{hyperref}
\usepackage[nameinlink,noabbrev]{cleveref}

\definecolor{TitleBlue}{HTML}{173A5E}
\definecolor{AccentBlue}{HTML}{2D668F}
\definecolor{SoftGray}{HTML}{F2F4F7}

\pagestyle{fancy}
\fancyhf{}
\fancyhead[L]{\small Surreal numbers as a real vector class}
\fancyhead[R]{\small\thepage}
\renewcommand{\headrulewidth}{0.3pt}

\setlength{\parindent}{1.15em}
\setlength{\parskip}{0.24em}
\setlength{\emergencystretch}{3em}
\widowpenalty=10000
\clubpenalty=10000
\displaywidowpenalty=10000
\allowdisplaybreaks[2]
\setlist{leftmargin=2.0em,itemsep=0.22em,topsep=0.35em}
\setcounter{tocdepth}{2}
\numberwithin{equation}{section}

\newtheorem{theorem}{Theorem}[section]
\newtheorem{proposition}[theorem]{Proposition}
\newtheorem{lemma}[theorem]{Lemma}
\newtheorem{corollary}[theorem]{Corollary}
\newtheorem{principle}[theorem]{Principle}
\theoremstyle{definition}
\newtheorem{definition}[theorem]{Definition}
\newtheorem{example}[theorem]{Example}
\newtheorem{question}[theorem]{Research question}
\theoremstyle{remark}
\newtheorem{remark}[theorem]{Remark}
\newtheorem{warning}[theorem]{Warning}

\crefname{theorem}{theorem}{theorems}
\crefname{proposition}{proposition}{propositions}
\crefname{lemma}{lemma}{lemmas}
\crefname{corollary}{corollary}{corollaries}
\crefname{definition}{definition}{definitions}
\crefname{example}{example}{examples}
\crefname{question}{research question}{research questions}
\crefname{warning}{warning}{warnings}
\crefname{principle}{principle}{principles}

\newcommand{\No}{\mathbf{No}}
\newcommand{\On}{\mathbf{On}}
\newcommand{\R}{\mathbb{R}}
\newcommand{\Q}{\mathbb{Q}}
\newcommand{\Z}{\mathbb{Z}}
\newcommand{\N}{\mathbb{N}}
\newcommand{\M}{\mathfrak{M}}
\newcommand{\Hahn}{\mathsf{H}}
\newcommand{\Fin}{\No_{\mathrm{fin}}}
\newcommand{\supp}{\operatorname{supp}}
\newcommand{\Span}{\operatorname{span}}
\newcommand{\lead}{\operatorname{lead}}
\newcommand{\coeff}{\operatorname{coeff}}
\newcommand{\gr}{\operatorname{gr}}
\newcommand{\Hom}{\operatorname{Hom}}
\newcommand{\End}{\operatorname{End}}
\newcommand{\Aut}{\operatorname{Aut}}
\newcommand{\trdeg}{\operatorname{trdeg}}
\newcommand{\id}{\operatorname{id}}
\newcommand{\ct}{\operatorname{ct}}
\newcommand{\im}{\operatorname{im}}
\newcommand{\Ord}{\operatorname{Ord}}
\newcommand{\sH}{\mathop{\sum}\nolimits^{\mathrm H}}
\newcommand{\abs}[1]{\lvert #1\rvert}
\newcommand{\norm}[1]{\lVert #1\rVert}
\newcommand{\angles}[1]{\langle #1\rangle}
\newcommand{\RepoCommit}{f7e7d8607d3b1c49c93b844cae20a77ab3689fe7}
\newcommand{\RepoShort}{f7e7d86}
\newcommand{\classprod}{\mathop{\prod}\nolimits^{\mathrm H}}

\hypersetup{
  pdftitle={The Surreal Numbers as a Real Vector Class},
  pdfauthor={Research report prepared for Vladimir Reshetnikov},
  pdfsubject={Conway--Hahn coordinates, Hamel bases, strong duality, and canonical operators},
  pdfkeywords={surreal numbers, vector space, Hahn basis, Hamel basis, Conway normal form, strong summation}
}

\newcolumntype{Y}{>{\raggedright\arraybackslash}X}
\newcolumntype{L}[1]{>{\raggedright\arraybackslash}p{#1}}

\begin{document}
\hypersetup{pageanchor=false}
\begin{titlepage}
\centering
\vspace*{-0.25in}
{\large\scshape A research report for the ProveIt project\par}
\vspace{0.22in}
{\Huge\bfseries\color{TitleBlue} The Surreal Numbers as a Real Vector Class\par}
\vspace{0.13in}
{\LARGE\color{AccentBlue} Conway--Hahn Coordinates, Hamel Bases,\par
Strong Duality, and Canonical Operators\par}
\vspace{0.22in}
{\large October 1, 2026\par}
\vspace{0.22in}
\begin{minipage}{0.94\textwidth}
\begin{abstract}
\noindent
The surreal numbers form a proper-class vector space over their canonical
copy of the real field.  The immediate question---what is its natural
basis?---has two different answers, depending on what ``basis'' means.
The canonical answer is the monomial class
\(
  \M=\{\omega^\gamma:\gamma\in\No\}.
\)
Every surreal has a unique Conway normal form, a set-supported,
reverse-well-ordered real linear combination of these monomials.  Thus
\(\M\) is a genuine \emph{Hahn basis}, a basis for strong transfinite
summation, and the canonical basis of the associated graded vector class.
It is not a Hamel basis, because Hamel combinations are finite and most
surreals have infinite normal-form support.

Working in NBG set theory with global choice, we prove that a Hamel class
basis does exist, but no set-sized family spans \(\No\), every such basis is
a proper class, and every Hamel basis contains a proper class of
infinite-support vectors.  The quotient by the finite-support skeleton
already contains an explicit proper-class linearly independent family.
On each infinite reverse-well-ordered support window \(A\), the set-sized
space of all series supported in \(A\) has Hamel dimension
\(\lvert\R\rvert^{\lvert A\rvert}\); in particular, a countable window has
continuum dimension, although its monomial Hahn basis is countable.

The normal form also yields a canonical Boolean algebra of support
projections and identifies the associated graded algebra with the real
group algebra of the additive surreal class.  We classify strongly real
linear maps by Hahn-compatible column families, and specialize this to an
exact classification of strongly summation-preserving real-valued
functionals.  Such a functional is a weighted coefficient sum whose active
exponents meet every reverse-well-ordered set only finitely; this permits
proper-class active sets, including the ordinal exponents.  In sharp
contrast, every order-bounded real linear functional is zero, hence there
is no nonzero positive real linear state and no nonzero order-compatible
real-valued seminorm on all of \(\No\).

Finally, the full natural valuation topology admits no non-eventually-
constant set-indexed convergent net, so Conway normal forms are strong
formal sums rather than Schauder expansions in that topology.  We describe
monomial-preserving ordered-field automorphisms, diagonal strong
derivations, and an explicit proper-class algebraically independent family.
A detailed research program closes the article.  Classical inputs are
Conway normal form, the omega map, Hahn--Neumann support arithmetic, and the
Erd\H{o}s--Kaplansky dimension theorem; the structural package assembled here
is offered as an unrefereed research draft, not as a claim of priority.
\end{abstract}
\end{minipage}
\vfill
\begin{minipage}{0.94\textwidth}\small
\textbf{Repository context.}  The report was prepared after inspecting
\texttt{VladimirReshetnikov/ProveIt}, especially its surreal-number reports
on vector and tensor fields, Hahn vector-space duality, infinite-dimensional
Hahn spectral theory, and automorphisms.  Repository snapshot:
\texttt{\RepoCommit} (\texttt{\RepoShort}).

\textbf{Status.}  This is an AI-assisted mathematical research draft with
complete proofs of the results stated as new deductions in this text.  It
has not been refereed and the new statements have not been formalized in
Lean.  Literature searches cannot certify novelty; the article therefore
separates classical inputs, elementary consequences, and proposed research
questions explicitly.
\end{minipage}
\end{titlepage}
\hypersetup{pageanchor=true}

\pagenumbering{roman}
\begingroup
\small
\setlength{\parskip}{0pt}
\tableofcontents
\endgroup
\clearpage
\pagenumbering{arabic}

\section{The question and the answer}\label{sec:intro}

The surreal numbers \(\No\) contain the real numbers as a canonical ordered
subfield.  Consequently their additive class is naturally a vector class over
\(\R\).  This observation is elementary; the word \emph{basis} is not.  In
ordinary linear algebra a basis means a Hamel basis, so every vector is a
\emph{finite} linear combination.  The defining coordinate theorem for
surreals, however, is Conway normal form, and its support may have any
set-sized ordinal order type.  Treating those two notions as interchangeable
obscures the main structure.

The answer developed in this article is:
\begin{quote}
\textbf{The canonical basis of \(\No\) over \(\R\) is the monomial class
\(\{\omega^\gamma:\gamma\in\No\}\) in the Hahn, strong-summation, and
associated-graded senses.  It is not a Hamel basis.  Hamel bases exist only
after a class-choice principle is specified, are necessarily proper classes,
and have no comparably canonical known description.}
\end{quote}

The distinction can be summarized as follows.

\begin{center}
\small
\begin{tabularx}{0.96\textwidth}{L{0.20\textwidth}Y Y}
\toprule
Notion of basis & Allowed expansion & Verdict for \(\No\) over \(\R\)\\
\midrule
Hamel basis & Finite linear combinations & Exists in NBG with global choice;
necessarily a proper class and noncanonical.\\
Conway--Hahn basis & Set-indexed, reverse-well-ordered, locally finite strong
sums & Canonically \(\{\omega^\gamma\}\); every surreal has one unique
expansion.\\
Associated-graded basis & One coordinate in each dominance grade &
Canonically the same monomials.\\
Schauder basis & Limits of partial sums in a topology & Not in the full fine
order or natural valuation topology: every set-indexed convergent net is
eventually constant.\\
Transcendence basis & Algebraic field generation & Choice-dependent; the
monomials nevertheless contain explicit proper-class algebraically
independent families.\\
\bottomrule
\end{tabularx}
\end{center}

\subsection{Why the repository makes this question timely}

The ProveIt repository already separates several uses of the word ``vector.''
Its report on vector and tensor fields studies finite tuples with surreal
coefficients.  Its reports on Hahn duality and Hahn spectral theory study
set-sized coefficient spaces \(V((t^\Gamma))\), strong sums, completion, and
several dual notions.  A separate report and the work of Kaplan, Krapp, and
Serra study automorphisms through Conway normal forms.  The present article
asks a more foundational question: view the scalar class \(\No\) itself as
an \(\R\)-vector object and determine which linear-algebraic structures are
canonical.

An informal MathOverflow discussion in 2013 already noticed both sides of
the issue: global choice should produce a class-sized real basis, while the
``homogeneous'' pieces ought to be real multiples of powers of \(\omega\)
\cite{MO2013}.  Conway normal form makes the latter intuition exact, but only
as a Hahn decomposition rather than an algebraic direct sum.

\subsection{Main results at a glance}\label{subsec:main-results}

Besides the central basis theorem, the article proves the following.
\begin{enumerate}[label=\textbf{\Alph*.}]
\item No set of surreals spans \(\No\) over \(\R\)
(\cref{thm:no-set-spans}).
\item Under global choice a Hamel class basis exists
(\cref{thm:class-basis}), but every one has a proper class of vectors outside
the finite-support skeleton (\cref{cor:basis-many-infinite}).
\item Every infinite reverse-well-ordered support window \(A\) has Hamel
dimension \(\lvert\R\rvert^{\lvert A\rvert}\)
(\cref{thm:window-dimension}).
\item Normal-form truncation gives canonical complementary projections for
every exponent subclass and a one-dimensional graded piece at every surreal
exponent (\cref{thm:support-splitting,thm:graded}).
\item Strongly linear maps are exactly Hahn-compatible column matrices
(\cref{thm:strong-matrix}), and the strong real dual is classified by
summation-finite coefficient weights (\cref{thm:strong-dual}).
\item The order-bounded real dual is zero
(\cref{thm:order-dual-zero}), although the strong dual is nonzero and even
contains functionals with proper-class active support.
\item The fine natural valuation topology has only eventually constant
set-indexed convergent nets (\cref{thm:eventually-constant}).
\item Monomial-preserving strong ordered-field automorphisms are controlled by
ordered exponent automorphisms and positive real characters
(\cref{thm:monomial-aut}), while additive real-valued characters of the
exponent group yield diagonal strong derivations
(\cref{thm:diagonal-derivation}).
\item Rational independence of exponents is exactly algebraic independence of
the corresponding monomials (\cref{thm:monomial-alg-ind}).
\end{enumerate}

\subsection{Claim status and novelty discipline}

The existence and uniqueness of Conway normal form, the omega map, Hahn
support arithmetic, and the Erd\H{o}s--Kaplansky theorem are classical.  Several
results below are immediate or short consequences once the correct basis
notion is fixed; others package familiar Hahn-space ideas in the proper-class
surreal setting.  The exact combination---especially the quotient witness,
strong real-dual classification, order-dual collapse, and the anti-Schauder
observation---was not located as a single published theorem package in the
targeted search.  This is not a certification that each statement is new.

\section{Foundations and Conway--Hahn coordinates}\label{sec:foundations}

\subsection{Class-theoretic convention}\label{subsec:classes}

We work in von Neumann--Bernays--Godel set theory with global choice (NBG+GC),
as is common in recent work on class-sized automorphism groups of \(\No\)
\cite{KKS}.  A \emph{class} may be a set or a proper class.  A class vector
space over the set field \(\R\) has the usual operations, with quantification
over its vectors understood class-theoretically.  A \emph{Hamel class basis}
is a linearly independent class such that every vector is a finite linear
combination of basis vectors.

The use of NBG+GC has three purposes.
\begin{enumerate}
\item It permits direct statements about \(\No\) as one algebraic object.
\item It permits a class-length recursion choosing compatible bases of its
set-sized birthday stages.
\item It keeps every actual infinite sum set-indexed.  No proper-class sum is
ever formed in this article.
\end{enumerate}
Readers preferring ZFC may interpret every construction inside a fixed
set-sized surreal subfield or birthday cutoff; all local arguments are then
ordinary set mathematics.  Only the global basis-existence statement and the
proper-class conclusions require the class framework.

\subsection{Reverse well ordering and supports}

\begin{definition}[Reverse well ordered]
A linearly ordered set \(S\) is \emph{reverse well ordered} if every nonempty
subset of \(S\) has a largest element.  Equivalently, \(S\) admits a unique
enumeration without repetition
\[
  S=\{\gamma_\beta:\beta<\lambda\},\qquad
  \gamma_0>\gamma_1>\cdots,
\]
for some ordinal \(\lambda\).
\end{definition}

Every subset of a reverse-well-ordered set is reverse well ordered.  This
simple hereditary property is responsible for the existence of arbitrary
support truncations later in the article.

\begin{theorem}[Conway normal form]\label{thm:cnf}
Every surreal number \(x\) has a unique expansion
\begin{equation}\label{eq:cnf}
  x=\sum_{\beta<\lambda} r_\beta\,\omega^{\gamma_\beta},
\end{equation}
where \(\lambda\) is an ordinal, the exponents
\(\gamma_\beta\in\No\) are strictly decreasing, and
\(r_\beta\in\R\setminus\{0\}\).  Conversely, every expression of this form
represents a surreal number.
\end{theorem}

This is the classical Conway--Gonshor normal-form theorem
\cite{Conway,Gonshor,vdDE,MantovaMatusinski}.  The map
\(\gamma\mapsto\omega^\gamma\) here is Conway's \emph{omega map}.  It sends
the additive ordered class \(\No\) isomorphically onto the multiplicative
ordered class of monomials and chooses the simplest positive representative
of each Archimedean equivalence class.  It should not be confused merely with
formal use of ordinary real exponentiation.

For \(x\) in the form \eqref{eq:cnf}, define
\[
  \supp(x)=\{\gamma_\beta:\beta<\lambda\},
  \qquad
  [\omega^\gamma]x=
  \begin{cases}
  r_\beta,&\gamma=\gamma_\beta,\\
  0,&\gamma\notin\supp(x).
  \end{cases}
\]
The first notation is the \emph{normal-form support}; the second is the
coefficient of \(\omega^\gamma\).  The support is always a set, even though
the exponent class is proper.

\begin{definition}[The Hahn product over the surreal exponent class]
Let
\begin{equation}\label{eq:hahn-product}
 \Hahn_{\R}(\No)=
 \left\{f:\No\to\R:
 \supp(f)=\{\gamma:f(\gamma)\ne0\}
 \text{ is a reverse-well-ordered set}\right\}.
\end{equation}
The functions are class functions, but each one has set support.  Addition and
real scalar multiplication are coefficientwise.
\end{definition}

The normal-form theorem gives a canonical \(\R\)-linear class isomorphism
\begin{equation}\label{eq:no-hahn-linear}
  \No\ \cong\ \Hahn_{\R}(\No),
  \qquad
  x\longmapsto\bigl(\gamma\mapsto[\omega^\gamma]x\bigr).
\end{equation}
Hahn--Neumann support arithmetic further identifies the ordered field
structure with the generalized power-series field
\begin{equation}\label{eq:no-hahn-field}
  \No=\R((\omega^{\No})).
\end{equation}
In particular,
\[
 \omega^\gamma\omega^\delta=\omega^{\gamma+\delta},
\]
and multiplication of two normal forms is the convolution of their
coefficients.  The support theorem guarantees that each output coefficient
is a finite sum and that the resulting support is reverse well ordered
\cite{Hahn,Neumann,Gonshor}.

\subsection{Strong Hahn summation}\label{subsec:strong-sum}

Normal form is not an ordinary finite sum and, as \cref{sec:topology} will
show, it is not generally a topological limit of finite partial sums.  The
appropriate operation is strong Hahn summation.

\begin{definition}[Hahn-summable family]\label{def:hahn-summable}
Let \(I\) be a set and \((x_i)_{i\in I}\) a family in \(\No\).  The family is
\emph{Hahn summable}, or \emph{strongly summable}, if
\begin{enumerate}
\item \(\bigcup_{i\in I}\supp(x_i)\) is reverse well ordered, and
\item for every exponent \(\gamma\), only finitely many \(i\) have
\([\omega^\gamma]x_i\ne0\).
\end{enumerate}
Its strong sum is
\begin{equation}\label{eq:strong-sum}
 \sH_{i\in I}x_i
 =\sum_{\gamma}
   \left(\sum_{i\in I}[\omega^\gamma]x_i\right)\omega^\gamma,
\end{equation}
where the inner sum is finite.
\end{definition}

This definition is purely algebraic and support-controlled.  It is invariant
under reindexing of the family and is associative under the usual flattening
hypotheses.  In particular, every normal form is the strong sum of its
monomial terms:
\begin{equation}\label{eq:normal-strong}
 x=\sH_{\gamma\in\supp(x)}
     [\omega^\gamma]x\,\omega^\gamma.
\end{equation}

\begin{example}[A geometric normal form]\label{ex:geometric}
The reverse-well-ordered set \(\{0,-1,-2,\ldots\}\) supports the identity
\begin{equation}\label{eq:geometric}
  \frac{1}{1-\omega^{-1}}
  =1+\omega^{-1}+\omega^{-2}+\cdots
  =\sH_{n<\omega}\omega^{-n}.
\end{equation}
The right side is a legitimate surreal normal form.  It is not a finite
linear combination of monomials.
\end{example}

\begin{example}[What is not allowed]
The class-indexed expression \(\sum_{\gamma\in\No}\omega^\gamma\) is not a
surreal number: its support is a proper class.  Likewise, a set-supported
formal expression whose support is not reverse well ordered is not a Conway
normal form.  The word ``transfinite'' therefore does not mean that arbitrary
class sums are permitted.
\end{example}

\subsection{Leading exponent, valuation, and Archimedean scale}

For \(x\ne0\), let
\[
 \lead(x)=\max\supp(x),
 \qquad
 \operatorname{lc}(x)=[\omega^{\lead(x)}]x.
\]
Then
\begin{equation}\label{eq:leading-factor}
 x=\operatorname{lc}(x)\,\omega^{\lead(x)}(1+\varepsilon),
 \qquad \abs{\varepsilon}<1/n\quad(n\in\N_{>0}),
\end{equation}
so \(\varepsilon\) is infinitesimal.  The sign of \(x\) is the sign of its
leading real coefficient, and the monomial \(\omega^{\lead(x)}\) is the
simplest positive representative of its Archimedean magnitude.

We use the additive natural valuation
\begin{equation}\label{eq:natural-valuation}
 v(x)=-\lead(x)\quad(x\ne0),\qquad v(0)=+\infty.
\end{equation}
Thus larger values mean smaller surreal magnitude.  This sign convention
agrees with writing the formal Hahn variable as \(t^\gamma=\omega^{-\gamma}\).

\section{The natural basis: Hahn, not Hamel}\label{sec:natural-basis}

\subsection{Four basis notions}

For a set-sized vector space, ``basis'' normally means Hamel basis.  In
non-Archimedean analysis and transseries, one also encounters topological,
orthogonal, and strong-summation bases.  The distinctions are essential here.

\begin{definition}[Hamel and Hahn bases]
Let \(V\) be a class vector space over \(\R\).
\begin{enumerate}
\item A \emph{Hamel class basis} \(B\) is linearly independent and every
\(x\in V\) is a unique finite sum \(\sum_{b\in F}r_b b\), with \(F\subset B\)
finite.
\item A \emph{Hahn basis} indexed by a linearly ordered class \(\Gamma\) is a
family \((e_\gamma)_{\gamma\in\Gamma}\) together with a specified strong
summation rule such that every vector has a unique expansion
\[
 x=\sH_{\gamma\in S}r_\gamma e_\gamma
\]
with \(S\) a set in the admissible support class.
\end{enumerate}
\end{definition}

A Hahn basis is not a weakened Hamel basis.  It is a basis for a richer
algebraic object: a vector space equipped with distinguished infinite sums.
Likewise, a Schauder basis requires a topology and convergence of partial
sums, a condition independent of formal Hahn summability.

\begin{theorem}[The canonical Conway--Hahn basis]\label{thm:hahn-basis}
The monomial class
\begin{equation}\label{eq:monomial-class}
  \M=\{\omega^\gamma:\gamma\in\No\}
\end{equation}
is a canonical Hahn basis of \(\No\) over \(\R\).  More precisely, the map
\[
 \Hahn_{\R}(\No)\longrightarrow\No,
 \qquad
 f\longmapsto\sH_{\gamma\in\supp(f)}f(\gamma)\omega^\gamma
\]
is an \(\R\)-linear class isomorphism, and its coordinate expansion is
unique.
\end{theorem}

\begin{proof}
Existence and uniqueness are exactly \cref{thm:cnf}.  Linearity follows from
coefficientwise addition and scalar multiplication.  Every support involved
is a reverse-well-ordered set, and no class-indexed sum is taken.
\end{proof}

The theorem answers the user's basis question in the sense best adapted to
surreal arithmetic.  It also explains why the answer is useful rather than
merely formal: the basis vectors multiply by addition of their indices, order
is read from the first nonzero coordinate, and every truncation of an
expansion is again a surreal.

\subsection{The finite-support skeleton}\label{subsec:finite-skeleton}

Define
\begin{equation}\label{eq:fin-def}
 \Fin=\{x\in\No:\supp(x)\text{ is finite}\}.
\end{equation}
Then
\begin{equation}\label{eq:direct-vs-hahn}
 \Fin=\bigoplus_{\gamma\in\No}\R\omega^\gamma
 \quad\subsetneq\quad
 \No=\classprod_{\gamma\in\No}\R\omega^\gamma.
\end{equation}
The symbol on the right denotes the Hahn product: set-supported,
reverse-well-ordered coordinate families, not the unrestricted direct
product.

\begin{proposition}[Where the monomials are a Hamel basis]
The class \(\M\) is a Hamel class basis of \(\Fin\), and it is real-linearly
independent in \(\No\).  It does not Hamel-span \(\No\).
\end{proposition}

\begin{proof}
A finite-support normal form is, by definition, a finite real linear
combination of distinct monomials, and uniqueness of normal form proves
linear independence.  Strictness follows from \cref{ex:geometric}, whose
support is infinite.
\end{proof}

Thus the most natural algebraic direct sum is not the whole surreal field.
Passing from \(\Fin\) to \(\No\) is analogous to passing from polynomials to a
Hahn-series field, not merely adjoining more basis vectors.

\begin{remark}[Why ``all Archimedean lines'' is not a direct-sum answer]
Each canonical line \(\R\omega^\gamma\) is Archimedean as an ordered additive
group, but \(\No\) is not their algebraic direct sum.  Moreover, there are
many nonhomogeneous one-dimensional real subspaces such as
\(\R(\omega+1)\).  The omega map is what canonically selects one pure line per
Archimedean scale.
\end{remark}

\subsection{Why the monomial basis is natural}\label{subsec:naturality}

The word ``natural'' can conceal taste.  Here it is supported by five precise
properties.

\begin{principle}[Simplicity normalization]
For every Archimedean class of positive surreals,
\(\omega^\gamma\) is its unique simplest positive representative.  Hence the
basis vectors are selected by the defining simplicity hierarchy of surreal
numbers, not by an external well ordering.
\end{principle}

\begin{principle}[Multiplicative grading]
The index law
\[
 \omega^\gamma\omega^\delta=\omega^{\gamma+\delta}
\]
turns the additive exponent class into the multiplicative monomial class.
The basis therefore interacts exactly with the field multiplication.
\end{principle}

\begin{principle}[Leading-coordinate order]
For nonzero \(x\), the sign and Archimedean scale of \(x\) are determined by
the coefficient and index of its largest supported monomial.  Thus the basis
is adapted to both the total order and the natural valuation.
\end{principle}

\begin{principle}[Truncation stability]
Every subset of a reverse-well-ordered support is reverse well ordered.
Consequently every coordinate truncation of a surreal normal form is again a
surreal number.
\end{principle}

\begin{principle}[Strong uniqueness]
The expansion remains unique even though arbitrary set-sized ordinal lengths
are allowed.  This is substantially stronger than a nonunique asymptotic
expansion.
\end{principle}

No arbitrary Hamel basis supplied by choice is known to enjoy this package.
Indeed, a Hamel basis must somehow encode infinite-support normal forms into
individual basis vectors, destroying the one-monomial-per-scale geometry.

\subsection{Changing the scalar field}\label{subsec:scalar-change}

The special role of \(\R\) can also be characterized by coefficient
normalization.

\begin{proposition}[Hahn bases over smaller coefficient fields]
Let \(k\subseteq\R\) be a subfield and let \(C\) be a Hamel basis of \(\R\)
over \(k\).  Then
\begin{equation}\label{eq:k-hahn-basis}
  \{c\omega^\gamma:c\in C,\ \gamma\in\No\}
\end{equation}
is a Hahn basis of \(\No\) over \(k\): every surreal has a unique strongly
summable expansion in these vectors, with finitely many \(C\)-coordinates at
each exponent.
\end{proposition}

\begin{proof}
Write each real normal-form coefficient \(r_\gamma\) uniquely as a finite
\(k\)-linear combination of elements of \(C\).  Replacing
\(r_\gamma\omega^\gamma\) by those finitely many terms preserves the union of
exponent supports and local finiteness.  Conversely, collecting the finitely
many terms at each exponent recovers the unique real coefficient and then the
unique Conway normal form.
\end{proof}

Over \(\R\), the coefficient basis is canonically \(\{1\}\).  Over \(\Q\),
for example, an additional highly noncanonical Hamel basis of \(\R\) is
unavoidable.  This is another sense in which \(\R\) is the natural scalar
field for the monomial coordinates.

\subsection{A first warning about topology}

It is tempting to enumerate the support in decreasing order and call the
normal form the limit of its finite initial sums.  In ordinary Hahn fields
this may work for selected supports and selected valuation topologies, but it
fails globally for \(\No\).  The full reason is proved in
\cref{sec:topology}; already the geometric series \eqref{eq:geometric} is not
the limit of its ordinary partial sums in the full natural valuation topology,
because the tail valuations \(1,2,3,\ldots\) never exceed the surreal
threshold \(\omega\).  Strong summation, not sequential convergence, is the
primitive operation.

\section{Hamel bases, class size, and local dimensions}\label{sec:hamel}

\subsection{No set-sized family can span}\label{subsec:no-set-span}

The first obstruction to ordinary linear algebra is stronger than
infinite-dimensionality.

\begin{theorem}[No set spans \(\No\)]\label{thm:no-set-spans}
If \(S\subseteq\No\) is a set, then
\[
  \Span_{\R}(S)\ne\No.
\]
More precisely, there is an exponent \(\gamma\in\No\) such that
\(\omega^\gamma\notin\Span_{\R}(S)\).
\end{theorem}

\begin{proof}
The union
\[
  U=\bigcup_{x\in S}\supp(x)
\]
is a set, by Replacement and Union.  The support of every finite real linear
combination of members of \(S\) is contained in \(U\), since addition and
scalar multiplication create no new exponents.  The class \(\No\) is proper,
so choose \(\gamma\notin U\).  Then \(\omega^\gamma\) has support
\(\{\gamma\}\) and cannot lie in the span.
\end{proof}

\begin{corollary}\label{cor:basis-proper}
Every Hamel basis of \(\No\) over \(\R\), if one is available in the chosen
class theory, is a proper class.  In particular, \(\No\) has no cardinal-valued
Hamel dimension in the ordinary set-theoretic sense.
\end{corollary}

The proof is support-theoretic and does not use birthdays, saturation, or
field multiplication.  It applies to many proper-class Hahn spaces.

\subsection{Existence with global choice}\label{subsec:basis-existence}

The ordinary vector-space basis theorem applies only to sets.  Nevertheless,
NBG+GC supplies a class analogue by a transparent stage construction.

For an ordinal \(\alpha\), let \(\No_{<\alpha}\) be the set of surreal
numbers born before day \(\alpha\), and put
\begin{equation}\label{eq:Valpha}
  V_\alpha=\Span_{\R}(\No_{<\alpha}).
\end{equation}
Each \(V_\alpha\) is a set-sized real vector space,
\(V_\alpha\subseteq V_\beta\) for \(\alpha<\beta\), and
\begin{equation}\label{eq:union-Va}
  \No=\bigcup_{\alpha\in\On}V_\alpha.
\end{equation}
At a limit ordinal \(\lambda\), every finite linear combination of elements
born below \(\lambda\) already occurs at some earlier stage, so
\(V_\lambda=\bigcup_{\alpha<\lambda}V_\alpha\).

\begin{theorem}[Existence of a Hamel class basis]\label{thm:class-basis}
In NBG with global choice, \(\No\) has a Hamel class basis over \(\R\).
\end{theorem}

\begin{proof}
Construct bases \(B_\alpha\) by transfinite class recursion.  Start with the
empty basis of \(V_0\).  At a successor stage, use the ordinary basis-extension
theorem inside the set \(V_{\alpha+1}\) to extend \(B_\alpha\) to a basis
\(B_{\alpha+1}\) of \(V_{\alpha+1}\).  Global choice permits these extensions
to be selected uniformly through all ordinals.  At a limit \(\lambda\), set
\(B_\lambda=\bigcup_{\alpha<\lambda}B_\alpha\); this is a basis of
\(V_\lambda\) because every vector and every finite dependence occurs at an
earlier stage.  Finally let \(B=\bigcup_{\alpha\in\On}B_\alpha\).  Every
surreal lies in some \(V_\alpha\) and has a finite \(B_\alpha\)-expansion,
while every finite dependence among members of \(B\) occurs in one stage.
Thus \(B\) is a Hamel class basis.
\end{proof}


\begin{proposition}[A Hamel basis may contain all monomials]
\label{prop:basis-containing-monomials}
In NBG with global choice, the linearly independent class \(\M\) extends to
a Hamel class basis \(B\) of \(\No\).  Every such extension has a proper
class of additional, infinite-support vectors.
\end{proposition}

\begin{proof}
Global choice gives a class enumeration
\((x_\alpha)_{\alpha\in\On}\) of \(\No\setminus\M\).  Recursively build a
set \(C_\alpha\subseteq\{x_\beta:\beta<\alpha\}\): at stage \(\alpha\), add
\(x_\alpha\) exactly when it is not in the finite real span of
\(\M\cup C_\alpha\).  Each stage adds at most one vector, so the previously
added part is a set and elementary transfinite recursion applies.  Put
\(C=\bigcup_{\alpha\in\On}C_\alpha\) and \(B=\M\cup C\).

A finite dependence in \(B\) would have a last-added member of \(C\), which
was added only because it was outside the span of all earlier members and
\(\M\); hence \(B\) is independent.  Every enumerated vector was either
added or already in the span at its stage, so \(B\) spans \(\No\).  Finally,
\cref{cor:basis-many-infinite} shows that \(B\setminus\Fin\), and therefore
also the supplement beyond the monomials, is a proper class.
\end{proof}

\begin{remark}[Existence is not naturality]
The construction makes class-many arbitrary extension choices.  It proves
that a Hamel basis exists under a stated foundation; it does not identify a
basis respected by normal-form truncations, multiplication, order, the omega
map, exponentiation, or the canonical derivations.  The Conway monomials are
canonical precisely because they solve a different, strong-summation basis
problem.
\end{remark}

\subsection{The quotient by finite support is already class-large}

The failure of the monomials to Hamel-span \(\No\) is not repaired by adding
only a set of infinite-support vectors.  We give an explicit witness.

For each ordinal \(\alpha\), define the disjoint reverse-well-ordered block
\begin{equation}\label{eq:block-Ea}
 E_\alpha=
 \{-\!(\omega\alpha+n):n<\omega\}\subseteq\No,
\end{equation}
where \(\omega\alpha\) denotes ordinal multiplication, viewed as a surreal
ordinal, and put
\begin{equation}\label{eq:block-ua}
 u_\alpha=\sH_{n<\omega}
          \omega^{-(\omega\alpha+n)}.
\end{equation}
The intervals \([\omega\alpha,\omega\alpha+\omega)\) are disjoint, hence so
are the supports \(E_\alpha\).

\begin{theorem}[A proper-class independent quotient family]
\label{thm:quotient-independent}
The cosets
\[
  \{u_\alpha+\Fin:\alpha\in\On\}
\]
are real-linearly independent in \(\No/\Fin\).
\end{theorem}

\begin{proof}
Take distinct ordinals \(\alpha_1,\ldots,\alpha_m\) and real numbers
\(c_1,\ldots,c_m\).  Because the supports \(E_{\alpha_j}\) are pairwise
disjoint, no cancellation occurs between different blocks.  If some
\(c_j\ne0\), the support of \(\sum_j c_j u_{\alpha_j}\) contains the entire
infinite block \(E_{\alpha_j}\), and therefore the sum does not belong to
\(\Fin\).  Hence a linear combination represents zero in the quotient only
when every coefficient is zero.
\end{proof}

\begin{corollary}\label{cor:basis-many-infinite}
Every Hamel class basis \(B\) of \(\No\) contains a proper class of vectors
outside \(\Fin\).
\end{corollary}

\begin{proof}
The images in \(\No/\Fin\) of the members of \(B\setminus\Fin\) span the
quotient.  If \(B\setminus\Fin\) were a set, its real span would be a set and
could not contain the proper-class family of distinct independent cosets from
\cref{thm:quotient-independent}.
\end{proof}

This result sharpens the slogan that a Hamel basis is nonmonomial.  It must
contain class-many genuinely infinite normal forms.

\subsection{Set-sized support windows}\label{subsec:windows}

Class size should not hide the familiar cardinal arithmetic visible in every
set-sized coordinate window.  Let \(A\subseteq\No\) be a reverse-well-ordered
set and define
\begin{equation}\label{eq:H-A}
  \Hahn(A)=\{x\in\No:\supp(x)\subseteq A\}.
\end{equation}
Every subset of \(A\) is reverse well ordered, so coefficient extraction gives
an ordinary vector-space isomorphism
\begin{equation}\label{eq:H-A-product}
  \Hahn(A)\cong\R^A.
\end{equation}
The monomials indexed by \(A\) form a Hahn basis of \(\Hahn(A)\), but only a
Hamel basis of its finite-support subspace.

\begin{theorem}[Hamel dimension of a support window]
\label{thm:window-dimension}
If \(A\) is an infinite reverse-well-ordered set, then
\begin{equation}\label{eq:window-dim}
  \dim_{\R}\Hahn(A)=\lvert\R\rvert^{\lvert A\rvert}.
\end{equation}
\end{theorem}

\begin{proof}
By \eqref{eq:H-A-product}, this is the Erd\H{o}s--Kaplansky theorem
\(\dim_D D^I=\lvert D^I\rvert\) for an infinite index set \(I\) and division
ring \(D\) \cite{Jacobson,Bergman}.  Taking \(D=\R\) and \(I=A\) gives the
formula.
\end{proof}

\begin{corollary}[A countable window has continuum dimension]
\label{cor:countable-window}
Let \(A=\{\gamma_0>\gamma_1>\cdots\}\) be countably infinite.  Then
\[
  \dim_{\R}\Hahn(A)=2^{\aleph_0}.
\]
An explicit continuum-sized independent family is
\begin{equation}\label{eq:vandermonde-family}
 x_t=\sH_{n<\omega}t^n\omega^{\gamma_n},
 \qquad t\in\R.
\end{equation}
\end{corollary}

\begin{proof}
The space has cardinality
\(\lvert\R\rvert^{\aleph_0}=2^{\aleph_0}\), so its dimension is at most the
continuum.  To prove the matching lower bound, take distinct
\(t_1,\ldots,t_m\in\R\) and suppose
\(\sum_{j=1}^m c_jx_{t_j}=0\).  Comparing the first \(m\) coordinates gives
\[
 \sum_{j=1}^m c_jt_j^n=0,
 \qquad n=0,\ldots,m-1.
\]
The Vandermonde matrix \((t_j^n)\) is invertible, hence all \(c_j=0\).
Thus \(\{x_t:t\in\R\}\) is linearly independent.
\end{proof}

\begin{example}[The gap between coordinate and Hamel dimensions]
On a countable support window, the canonical monomial Hahn basis has
cardinality \(\aleph_0\), while every Hamel basis has cardinality continuum.
The difference is exactly the information carried by arbitrary infinite
coordinate strings.
\end{example}

\subsection{What a ``dimension of \texorpdfstring{\(\No\)}{No}'' can mean}

There is no set cardinal \(\kappa\) with \(\dim_{\R}\No=\kappa\), because no
set spans.  Three replacements are meaningful.
\begin{enumerate}
\item The \emph{class Hamel dimension} is represented by any Hamel basis
class, but comparing proper-class sizes requires the chosen class theory and
is too coarse to capture support.
\item The \emph{graded dimension} is one at each exponent
(\cref{sec:filtration}).
\item The \emph{window dimension profile}
\(A\mapsto\lvert\R\rvert^{\lvert A\rvert}\) records the ordinary Hamel
dimensions of set-sized Hahn blocks.
\end{enumerate}
The second and third are canonical and structurally informative; the first is
mostly an existence invariant.

\section{Canonical projections, flags, and the associated graded}\label{sec:filtration}

The monomial coordinates provide much more than a list of generators.  They
split \(\No\) functorially by arbitrary exponent classes and resolve its
natural valuation into one-dimensional layers.

\subsection{Support projections}\label{subsec:support-projections}

For a subclass \(A\subseteq\No\), define
\begin{equation}\label{eq:PA}
 P_A(x)=\sH_{\gamma\in\supp(x)\cap A}
       [\omega^\gamma]x\,\omega^\gamma.
\end{equation}
This is well defined because \(\supp(x)\cap A\) is a subset of a
reverse-well-ordered set.

\begin{theorem}[Canonical support splitting]\label{thm:support-splitting}
For every subclass \(A\subseteq\No\), \(P_A\) is a strongly summation-
preserving real-linear idempotent and
\begin{equation}\label{eq:support-split}
 \No=\No[A]\oplus\No[\No\setminus A],
 \qquad
 \No[A]=\{x:\supp(x)\subseteq A\}.
\end{equation}
The projections satisfy
\begin{align}
 P_AP_B&=P_{A\cap B},\label{eq:boolean1}\\
 P_{\No\setminus A}&=\id-P_A,\label{eq:boolean2}\\
 P_{A\cup B}&=P_A+P_B-P_{A\cap B}.\label{eq:boolean3}
\end{align}
Thus exponent subclasses act as a Boolean algebra of canonical linear
projections.
\end{theorem}

\begin{proof}
Coefficientwise linearity and idempotence are immediate.  Every normal form
splits uniquely into the terms with exponent in \(A\) and those outside
\(A\), proving \eqref{eq:support-split}.  If \((x_i)\) is Hahn summable, the
union of the supports of \((P_Ax_i)\) is a subset of the original
reverse-well-ordered union, and local finiteness at every exponent is
preserved.  Hence \(P_A\) is strong.  The Boolean identities are checked on
each coordinate.
\end{proof}

Important special cases include:
\begin{enumerate}
\item the single-coordinate projection
\(P_{\{\gamma\}}x=[\omega^\gamma]x\,\omega^\gamma\);
\item initial and final truncations at an exponent \(\delta\);
\item the decomposition into positive, zero, and negative exponents;
\item projections onto ordinal exponents, real exponents, or any definable
scale family.
\end{enumerate}
No Hamel-basis choice is needed.

\begin{warning}
Support projections are real-linear, not generally field homomorphisms.  For
example, constant-term extraction is multiplicative on certain support
subrings but not on all of \(\No\), because exponents from different factors
can add to zero.
\end{warning}

\subsection{The dominance flag}

For \(\gamma\in\No\), define
\begin{align}
 F_{\le\gamma}
 &=\{x\in\No:\supp(x)\subseteq(-\infty,\gamma]\},\label{eq:Fle}\\
 F_{<\gamma}
 &=\{x\in\No:\supp(x)\subseteq(-\infty,\gamma)\}.
 \label{eq:Fl}
\end{align}
Equivalently, a nonzero \(x\) belongs to \(F_{\le\gamma}\) exactly when
\(\lead(x)\le\gamma\).  These subspaces form an increasing proper-class flag.

\begin{proposition}[One-dimensional graded quotients]\label{prop:grade-one}
For every \(\gamma\in\No\), coefficient extraction induces a canonical
isomorphism
\begin{equation}\label{eq:one-grade}
  F_{\le\gamma}/F_{<\gamma}\ \cong\ \R,
  \qquad
  x+F_{<\gamma}\longmapsto[\omega^\gamma]x.
\end{equation}
The class of \(\omega^\gamma\) is the distinguished basis vector of this
quotient.
\end{proposition}

\begin{proof}
The kernel of the coefficient map on \(F_{\le\gamma}\) is precisely
\(F_{<\gamma}\), and every real coefficient occurs as
\(r\omega^\gamma\).
\end{proof}

Thus \(\No\) has one real direction at every Archimedean scale, plus the
support-completion needed to combine set-many scales.

\subsection{Associated graded algebra}\label{subsec:graded}

Let
\begin{equation}\label{eq:gr-def}
  \gr(\No)=\bigoplus_{\gamma\in\No}
  F_{\le\gamma}/F_{<\gamma}.
\end{equation}
The direct sum is algebraic: an element has only finitely many nonzero
homogeneous coordinates.  Multiplication of leading terms makes this a graded
\(\R\)-algebra.

\begin{theorem}[The associated graded is the group algebra]
\label{thm:graded}
There is a canonical graded \(\R\)-algebra isomorphism
\begin{equation}\label{eq:gr-group-algebra}
  \gr(\No)\ \cong\ \R[\No_{\mathrm{add}}]
  \ \cong\ \Fin,
\end{equation}
where \(\R[\No_{\mathrm{add}}]\) is the finite-support group algebra of the
additive surreal class.  Under this isomorphism the homogeneous symbol of
\(\omega^\gamma\) is the group-algebra element \(e_\gamma\), with
\[
 e_\gamma e_\delta=e_{\gamma+\delta}.
\]
\end{theorem}

\begin{proof}
By \cref{prop:grade-one}, the grade at \(\gamma\) is one-dimensional with
basis the symbol of \(\omega^\gamma\).  If
\(x=r\omega^\gamma+\text{lower terms}\) and
\(y=s\omega^\delta+\text{lower terms}\), then
\[
 xy=rs\omega^{\gamma+\delta}+\text{terms of exponent below }\gamma+\delta.
\]
Hence multiplication of homogeneous symbols is exactly the group-algebra
law.  The finite-support normal forms identify that group algebra with
\(\Fin\).
\end{proof}

\begin{remark}[What completion adds]
The associated graded remembers every individual dominance layer and their
multiplicative law, but it forgets the admissible infinite assembly of those
layers.  Recovering \(\No\) from \(\gr(\No)\) therefore requires the
reverse-well-ordered support rule.  This is the precise sense in which
\(\No\) is a Hahn completion of its graded skeleton, rather than an ordinary
linear span.
\end{remark}

\subsection{Canonical coefficient exact sequences}

For any exponent subclass \(A\), restriction of coefficients gives a split
exact sequence of real vector classes
\begin{equation}\label{eq:split-exact-A}
  0\longrightarrow \No[\No\setminus A]
  \longrightarrow \No
  \xrightarrow{\ P_A\ } \No[A]
  \longrightarrow0.
\end{equation}
For a singleton \(A=\{\gamma\}\), composition with
\(r\omega^\gamma\mapsto r\) gives
\begin{equation}\label{eq:coord-exact}
  0\longrightarrow\ker[\omega^\gamma]
  \longrightarrow\No
  \xrightarrow{[\omega^\gamma]}\R
  \longrightarrow0.
\end{equation}
These canonical splittings will supply the simplest nonzero members of the
strong dual in \cref{sec:strong-maps}.

\section{Strong linear maps and the intrinsic real dual}\label{sec:strong-maps}

A Hamel-linear map is determined by a Hamel basis, hence inherits all of its
choice dependence.  A map adapted to Conway normal form should instead
commute with the distinguished strong sums.

\subsection{Strong linearity}

\begin{definition}[Strong real-linear map]\label{def:strong-map}
An \(\R\)-linear class map \(T:\No\to\No\) is \emph{strong} if for every
Hahn-summable set-indexed family \((x_i)_{i\in I}\), the family
\((Tx_i)_{i\in I}\) is Hahn summable and
\begin{equation}\label{eq:strong-map}
  T\!\left(\sH_{i\in I}x_i\right)=\sH_{i\in I}Tx_i.
\end{equation}
\end{definition}

The definition includes both preservation of admissibility and commutation
with the sum.  Merely writing a coefficientwise formula is not enough unless
its output supports satisfy the Hahn conditions.

\begin{definition}[Hahn-compatible column family]
A class family \((y_\gamma)_{\gamma\in\No}\) in \(\No\) is
\emph{Hahn compatible} if for every reverse-well-ordered set
\(S\subseteq\No\), the family \((y_\gamma)_{\gamma\in S}\) is Hahn summable.
\end{definition}

Write
\begin{equation}\label{eq:matrix-entry}
  y_\gamma=\sH_{\delta}
  a_{\delta\gamma}\omega^\delta,
  \qquad a_{\delta\gamma}\in\R.
\end{equation}
Then compatibility says exactly that, for every reverse-well-ordered set
\(S\),
\begin{enumerate}
\item \(\bigcup_{\gamma\in S}\{\delta:a_{\delta\gamma}\ne0\}\) is reverse
well ordered, and
\item for each row exponent \(\delta\), only finitely many
\(\gamma\in S\) have \(a_{\delta\gamma}\ne0\).
\end{enumerate}
This is the surreal-Hahn analogue of a column-finiteness condition.

\begin{theorem}[Hahn-matrix classification of strong maps]
\label{thm:strong-matrix}
The assignments
\begin{equation}\label{eq:column-correspondence}
  T\longmapsto(T\omega^\gamma)_{\gamma\in\No}
\end{equation}
and
\begin{equation}\label{eq:T-from-columns}
  (y_\gamma)\longmapsto
  \left[
  \sH_{\gamma\in S}r_\gamma\omega^\gamma
  \longmapsto
  \sH_{\gamma\in S}r_\gamma y_\gamma
  \right]
\end{equation}
are inverse bijections between strong \(\R\)-linear endomorphisms of \(\No\)
and Hahn-compatible column families.
\end{theorem}

\begin{proof}
Let \(T\) be strong.  For a reverse-well-ordered set \(S\), the monomial
family \((\omega^\gamma)_{\gamma\in S}\) is Hahn summable.  Therefore
\((T\omega^\gamma)_{\gamma\in S}\) is Hahn summable, so the columns are
compatible.  Applying strong linearity to a normal form gives
\eqref{eq:T-from-columns}, establishing uniqueness.

Conversely, suppose \((y_\gamma)\) is compatible and define \(T\) by
\eqref{eq:T-from-columns}.  The output sum exists by compatibility and is
plainly real linear.  Let \((x_i)_{i\in I}\) be Hahn summable and write
\(x_i=\sH_\gamma r_{i\gamma}\omega^\gamma\).  Put
\(S=\bigcup_i\supp(x_i)\).  Compatibility says that
\((y_\gamma)_{\gamma\in S}\) has reverse-well-ordered union of supports and
is locally finite in every output exponent.  For a fixed \(\gamma\), only
finitely many \(i\) have \(r_{i\gamma}\ne0\).  Hence the flattened family
\((r_{i\gamma}y_\gamma)_{(i,\gamma)}\) is Hahn summable: at each output
exponent only finitely many \(\gamma\)'s contribute, and for each such
\(\gamma\) only finitely many \(i\)'s contribute.  Associativity of finite
coefficient sums now yields
\[
 T\!\left(\sH_i x_i\right)
 =\sH_{i,\gamma}r_{i\gamma}y_\gamma
 =\sH_i Tx_i.
\]
Thus \(T\) is strong.
\end{proof}

\begin{remark}
The theorem is a universal property of the Conway--Hahn basis.  A Hamel basis
determines all algebraic linear maps; the monomial Hahn basis determines all
strong linear maps, with the compatibility condition recording exactly which
images of basis vectors extend to every normal form.
\end{remark}

\subsection{Diagonal maps and support projections}

For an arbitrary class function \(a:\No\to\R\), define
\begin{equation}\label{eq:Da}
  D_a\!\left(\sH_\gamma r_\gamma\omega^\gamma\right)
  =\sH_\gamma a(\gamma)r_\gamma\omega^\gamma.
\end{equation}
The output support is contained in the input support, so \(D_a\) is strong.

\begin{proposition}[Diagonal operator calculus]\label{prop:diagonal}
For class functions \(a,b:\No\to\R\),
\begin{align}
 D_aD_b&=D_{ab},\\
 D_a+D_b&=D_{a+b}.
\end{align}
Moreover:
\begin{enumerate}
\item \(D_a\) is injective exactly when \(a(\gamma)\ne0\) for every
\(\gamma\);
\item \(D_a\) is a strong automorphism exactly when every \(a(\gamma)\ne0\),
with inverse \(D_{1/a}\);
\item \(D_a\) is idempotent exactly when
\(a(\gamma)\in\{0,1\}\) for every \(\gamma\).
\end{enumerate}
In particular, the support projection \(P_A\) is
\(D_{\mathbf{1}_A}\).
\end{proposition}

\begin{proof}
All assertions follow by checking each monomial coordinate.  Over the field
\(\R\), the scalar equation \(a^2=a\) has only the solutions \(0\) and
\(1\).
\end{proof}

\subsection{The strong real-valued dual}\label{subsec:strong-dual}

Embed \(\R\) in \(\No\) as the exponent-zero line.  A family of real numbers
is Hahn summable in this subspace exactly when only finitely many members are
nonzero, because all nonzero terms occupy the same exponent \(0\).

\begin{definition}[Summation-finite exponent class]\label{def:summation-finite}
A subclass \(A\subseteq\No\) is \emph{summation finite} if
\begin{equation}\label{eq:summation-finite}
  A\cap S\text{ is finite for every reverse-well-ordered set }
  S\subseteq\No.
\end{equation}
A weight \(a:\No\to\R\) is summation finite if its active class
\(\supp(a)=\{\gamma:a(\gamma)\ne0\}\) is summation finite.
\end{definition}


\begin{proposition}[Sequential characterization]
\label{prop:summation-finite-seq}
A subclass \(A\subseteq\No\) is summation finite if and only if it contains
no strictly decreasing sequence
\(a_0>a_1>a_2>\cdots\) indexed by \(\omega\).
\end{proposition}

\begin{proof}
A strictly decreasing sequence has reverse-well-ordered range: every nonempty
subset of its range has a largest member, corresponding to the least index
that occurs.  Hence such a sequence violates summation finiteness.
Conversely, if \(A\cap S\) is infinite for a reverse-well-ordered set \(S\),
enumerate that intersection in decreasing order.  Its order type is an
infinite ordinal and its first \(\omega\) members form a strictly decreasing
sequence in \(A\).
\end{proof}

\begin{theorem}[Classification of the strong real dual]
\label{thm:strong-dual}
An \(\R\)-linear map \(\Lambda:\No\to\R\) is strong if and only if there is a
summation-finite weight \(a:\No\to\R\) such that
\begin{equation}\label{eq:strong-functional}
  \Lambda(x)=
  \sum_{\gamma\in\supp(x)\cap\supp(a)}
  a(\gamma)[\omega^\gamma]x.
\end{equation}
The sum in \eqref{eq:strong-functional} is always finite.  The weight is
unique and equals
\begin{equation}\label{eq:weight-from-Lambda}
  a(\gamma)=\Lambda(\omega^\gamma).
\end{equation}
\end{theorem}

\begin{proof}
Suppose \(\Lambda\) is strong and define \(a\) by
\eqref{eq:weight-from-Lambda}.  For every reverse-well-ordered set \(S\), the
family \((\omega^\gamma)_{\gamma\in S}\) is Hahn summable.  Its image family
\((a(\gamma))_{\gamma\in S}\) must be summable in \(\R\), hence only
finitely many of those reals are nonzero.  Thus \(\supp(a)\) is summation
finite.  Applying \(\Lambda\) to the strong normal form of \(x\) yields
\eqref{eq:strong-functional}.

Conversely, let \(a\) be summation finite and define \(\Lambda\) by
\eqref{eq:strong-functional}.  It is real linear.  If \((x_i)_{i\in I}\) is
Hahn summable and \(S=\bigcup_i\supp(x_i)\), then
\(F=S\cap\supp(a)\) is finite.  For each \(\gamma\in F\), only finitely many
\(i\) have a nonzero \(\gamma\)-coefficient.  Therefore only finitely many
\(\Lambda(x_i)\) can be nonzero, and finite rearrangement of coefficients
gives
\(\Lambda(\sH_i x_i)=\sum_i\Lambda(x_i)\).  Hence \(\Lambda\) is strong.
Uniqueness follows by evaluation on monomials.
\end{proof}

\begin{example}[Coordinate functionals]
For every \(\gamma\), the coefficient map
\[
  \pi_\gamma(x)=[\omega^\gamma]x
\]
is strong.  It corresponds to the weight supported on the singleton
\(\{\gamma\}\).
\end{example}

The next example shows that a strong functional need not use only finitely
many coordinates globally.

\begin{lemma}[Well-ordered active classes are summation finite]
\label{lem:well-ordered-active}
If a subclass \(A\subseteq\No\) is well ordered by the inherited order, then
it is summation finite.
\end{lemma}

\begin{proof}
For a reverse-well-ordered set \(S\), the intersection \(A\cap S\) is both
well ordered and reverse well ordered.  An infinite linear order cannot have
both properties: an infinite ordinal, or any infinite set of ordinal type,
has a subset with no greatest element.  Hence \(A\cap S\) is finite.
\end{proof}

\begin{example}[The ordinal-coefficient trace]\label{ex:ordinal-trace}
The class \(\On\) is well ordered, so
\begin{equation}\label{eq:ordinal-trace}
  \operatorname{Tr}_{\On}(x)
  =\sum_{\gamma\in\supp(x)\cap\On}
    [\omega^\gamma]x
\end{equation}
is a well-defined strong \(\R\)-linear functional.  Its active exponent
class is the proper class \(\On\).  Every individual normal form contains
only finitely many ordinal exponents, even though the functional is active at
all ordinals.
\end{example}


\begin{remark}[Three duals already appear]
There are at least three very different notions of real dual for \(\No\):
\begin{enumerate}
\item the algebraic dual determined by a choice of Hamel basis;
\item the intrinsic strong dual classified by \cref{thm:strong-dual};
\item the order-bounded dual, which is zero by \cref{thm:order-dual-zero}.
\end{enumerate}
This mirrors, at the scalar-class level, the repository's separation of
represented, strong, and valuation-continuous duals for set-sized Hahn vector
spaces.
\end{remark}

\section{Order duality and the failure of Schauder convergence}\label{sec:topology}

The strong dual is nontrivial because it is adapted to formal support sums.
Demanding compatibility with the total order has the opposite effect: every
real-valued order-bounded linear functional vanishes.

\subsection{The order-bounded real dual is zero}

\begin{definition}[Order-bounded linear map]
An \(\R\)-linear map \(L:\No\to\R\) is \emph{order bounded} if for every
\(y\ge0\) there is a real number \(M_y\) such that
\[
  \abs{x}\le y\quad\Longrightarrow\quad\abs{L(x)}\le M_y.
\]
Equivalently, it maps every surreal order interval to a bounded subset of
\(\R\).
\end{definition}

\begin{theorem}[Collapse of the order-bounded real dual]
\label{thm:order-dual-zero}
Every order-bounded \(\R\)-linear map \(L:\No\to\R\) is zero.
\end{theorem}

\begin{proof}
Fix \(x\in\No\).  If \(x=0\) there is nothing to prove.  Put
\(y=\omega\abs{x}>0\).  For every positive integer \(n\), one has
\(n\abs{x}<\omega\abs{x}=y\), so \(nx\in[-y,y]\).  Order boundedness gives a
real \(M_y\) with
\[
  n\abs{L(x)}=\abs{L(nx)}\le M_y
  \qquad(n\in\N_{>0}).
\]
No positive real number can satisfy these inequalities for every \(n\).
Hence \(L(x)=0\).  Since \(x\) was arbitrary, \(L=0\).
\end{proof}

\begin{corollary}[No positive real linear state]\label{cor:no-positive-state}
Every positive \(\R\)-linear map \(L:\No\to\R\) is zero.  In particular,
there is no positive real-linear retraction of \(\No\) onto its real
subfield.
\end{corollary}

\begin{proof}
If \(L\) is positive and \(-y\le x\le y\), then
\(-L(y)\le L(x)\le L(y)\).  Thus \(L\) is order bounded and
\cref{thm:order-dual-zero} applies.
\end{proof}

\begin{example}[Constant term is strong but not positive]
The coefficient functional
\[
  \ct(x)=[\omega^0]x
\]
is strong and restricts to the identity on \(\R\), but it is not positive:
\(\omega-1>0\) while \(\ct(\omega-1)=-1\).  This example shows exactly why
\cref{cor:no-positive-state} does not contradict the existence of canonical
coefficient retractions.
\end{example}

\begin{corollary}[No nonzero order-compatible real seminorm]
\label{cor:no-order-seminorm}
Suppose \(p:\No\to\R_{\ge0}\) is absolutely homogeneous over \(\R\) and
monotone with respect to surreal magnitude:
\[
  \abs{x}\le\abs{y}\quad\Longrightarrow\quad p(x)\le p(y).
\]
Then \(p=0\).
\end{corollary}

\begin{proof}
For \(x\ne0\) and every \(n\in\N_{>0}\),
\(\abs{nx}<\abs{\omega x}\).  Hence
\[
 n p(x)=p(nx)\le p(\omega x),
\]
where the right side is one fixed real number.  Therefore \(p(x)=0\).
\end{proof}

A choice of Hamel basis can of course be declared orthonormal, producing a
noncanonical real-valued norm.  The corollary says that no such norm can
respect the ambient surreal order in the displayed sense.

\subsection{The fine valuation topology}

The natural valuation \eqref{eq:natural-valuation} gives fine balls
\begin{equation}\label{eq:valuation-ball}
  B_\gamma(0)=\{x\in\No:v(x)>\gamma\},
  \qquad \gamma\in\No.
\end{equation}
Translates of these balls define the full natural valuation topology.  It is
``fine'' because every surreal threshold is allowed, not merely thresholds
from a fixed set-sized subgroup.

The following elementary fact about \(\No\) is decisive.

\begin{lemma}[Every set of surreals is bounded]
\label{lem:set-bounded}
Every set \(A\subseteq\No\) has a surreal upper bound and a surreal lower
bound.
\end{lemma}

\begin{proof}
The Conway cut \(\{A\mid\varnothing\}\) is greater than every member of
\(A\).  Applying this to \(-A\) gives a lower bound.
\end{proof}

\begin{theorem}[Set-indexed convergence is eventually constant]
\label{thm:eventually-constant}
Let \((x_i)_{i\in I}\) be a net in \(\No\) whose directed index \(I\) is a
set.  If \(x_i\to x\) in the full natural valuation topology, then
\(x_i=x\) eventually.
\end{theorem}

\begin{proof}
Consider the set
\[
 A=\{v(x_i-x):i\in I,\ x_i\ne x\}\subseteq\No.
\]
By \cref{lem:set-bounded}, choose \(\beta\in\No\) strictly above every member
of \(A\).  Convergence to \(x\) in the neighborhood
\(x+B_\beta(0)\) implies that eventually
\(v(x_i-x)>\beta\).  No nonzero difference in the net has such a value by
the choice of \(\beta\), so eventually \(x_i-x=0\).
\end{proof}

\begin{remark}[Non-discrete but net-rigid]
The topology is not discrete: no single threshold \(\gamma\) excludes every
nonzero infinitesimal, because \(\omega^{-\delta}\in B_\gamma(0)\) for
sufficiently large \(\delta\).  Nevertheless, every \emph{set}-indexed net
has only a set of valuation values, hence a common upper bound.  This is a
proper-class phenomenon: the neighborhood system itself has proper-class
cofinality.
\end{remark}

\begin{corollary}[Normal forms are not Schauder expansions]
\label{cor:not-schauder}
If a surreal has infinite normal-form support, no set-indexed net of its
finite support truncations converges to it in the full natural valuation
topology.  Consequently the Conway monomials are not a Schauder basis for
that topology.
\end{corollary}

\begin{proof}
A convergent net of finite truncations would be eventually equal to its
limit by \cref{thm:eventually-constant}, forcing the limit to have finite
support.
\end{proof}

The same argument applies to the fine order topology.

\begin{corollary}[Fine order convergence is also eventually constant]
\label{cor:order-net}
Every set-indexed net convergent in the order topology generated by all
surreal intervals \((x-\varepsilon,x+\varepsilon)\), with
\(\varepsilon>0\), is eventually constant.
\end{corollary}

\begin{proof}
For the set of nonzero distances
\(D=\{\abs{x_i-x}:x_i\ne x\}\), the cut \(\{0\mid D\}\) gives a positive
surreal \(\varepsilon\) smaller than every member of \(D\).  Eventual
membership in the \(\varepsilon\)-interval around \(x\) therefore forces
equality.
\end{proof}

\subsection{How nontrivial convergence can be recovered}

The anti-Schauder theorem concerns the \emph{full} value group \(\No\).
Useful analytic work generally chooses a set-sized Hahn workspace
\(k((t^\Gamma))\), a set-sized cofinal scale family, or a coarser topology.
Within such a workspace, a support can be cofinal in the chosen value group,
and partial sums may converge.  This is why the repository's analytical
reports carefully distinguish set-sized workspaces, strong summation, and
fine surreal topology.  A future theory of ``Schauder bases for \(\No\)''
must specify which scale system has been retained and which proper-class
neighborhoods have been discarded.

\section{Monomial symmetries, derivations, and algebraic independence}
\label{sec:multiplicative}

The Hahn basis is not only linear.  Because its indices form the additive
surreal group, linear transformations of the monomials often extend to field
automorphisms and derivations.

\subsection{Monomial-preserving ordered-field automorphisms}

Let \(\phi:(\No,+,<)\to(\No,+,<)\) be an ordered additive class
automorphism, and let
\[
  \chi:(\No,+)\longrightarrow(\R_{>0},\cdot)
\]
be a positive real character.  Define
\begin{equation}\label{eq:T-chiphi}
 T_{\chi,\phi}\!\left(
   \sH_{\gamma\in S}r_\gamma\omega^\gamma\right)
 =\sH_{\gamma\in S}
   r_\gamma\chi(\gamma)\omega^{\phi(\gamma)}.
\end{equation}

\begin{theorem}[Monomial automorphisms]\label{thm:monomial-aut}
The map \(T_{\chi,\phi}\) is a strong \(\R\)-linear ordered-field
automorphism of \(\No\).  Conversely, every strong \(\R\)-linear
ordered-field automorphism that permutes the monomial lines
\(\{\R\omega^\gamma:\gamma\in\No\}\) has this form for unique
\(\chi\) and \(\phi\).
\end{theorem}

\begin{proof}
Because \(\phi\) is an order isomorphism, it sends reverse-well-ordered sets
to reverse-well-ordered sets.  The columns in \eqref{eq:T-chiphi} have
singleton supports with distinct output exponents, so
\cref{thm:strong-matrix} gives strong real linearity.  The inverse is obtained
from \(\phi^{-1}\) and the character
\(\gamma\mapsto\chi(\phi^{-1}\gamma)^{-1}\).

For monomials,
\begin{align*}
 T_{\chi,\phi}(\omega^\gamma\omega^\delta)
 &=\chi(\gamma+\delta)\omega^{\phi(\gamma+\delta)}\\
 &=\chi(\gamma)\chi(\delta)
   \omega^{\phi(\gamma)+\phi(\delta)}\\
 &=T_{\chi,\phi}(\omega^\gamma)
   T_{\chi,\phi}(\omega^\delta).
\end{align*}
Strong distributivity extends this identity to arbitrary normal forms.
Positive real coefficients and the increasing map \(\phi\) preserve the sign
of the leading term, hence the order.

Conversely, suppose \(T\) strongly permutes monomial lines.  Write uniquely
\[
  T(\omega^\gamma)=\chi(\gamma)\omega^{\phi(\gamma)},
  \qquad \chi(\gamma)>0.
\]
Bijectivity on lines makes \(\phi\) a bijection.  Multiplicativity of \(T\)
gives
\(\phi(\gamma+\delta)=\phi(\gamma)+\phi(\delta)\) and
\(\chi(\gamma+\delta)=\chi(\gamma)\chi(\delta)\).  Preservation of the order
of positive monomials makes \(\phi\) increasing.  Strong linearity then
forces \(T\) to act on every normal form by \eqref{eq:T-chiphi}.
\end{proof}

This elementary subgroup is one of the reasons the monomial basis is useful:
its symmetries are controlled by familiar group data.  The full automorphism
group is richer; Kaplan, Krapp, and Serra give a systematic Hahn-field
decomposition and distinguish strongly linear automorphisms from general
ones \cite{KKS}.

\begin{example}[Exponent dilations]\label{ex:exponent-dilation}
For every positive surreal \(a\), multiplication
\[
  \phi_a(\gamma)=a\gamma
\]
is an ordered additive automorphism of \(\No\).  Taking \(\chi=1\) gives
\begin{equation}\label{eq:sigma-a}
  \sigma_a\!\left(\sH_\gamma r_\gamma\omega^\gamma\right)
  =\sH_\gamma r_\gamma\omega^{a\gamma}.
\end{equation}
Then \(\sigma_a(\omega)=\omega^a\),
\(\sigma_a^{-1}=\sigma_{a^{-1}}\), and
\(\sigma_a\sigma_b=\sigma_{ab}\).
These are ordered-field automorphisms fixing \(\R\).  Preservation of the
surreal exponential or of the omega map as an additional named operation is
a separate and much more rigid requirement.
\end{example}

\begin{theorem}[Dilation-invariant strong functionals]
\label{thm:dilation-invariant-dual}
Let \(\sigma_a\) be the exponent-dilation automorphism
\eqref{eq:sigma-a}.  A strong real-linear functional
\(\Lambda:\No\to\R\) satisfies
\[
  \Lambda\circ\sigma_a=\Lambda
  \qquad\text{for every }a\in\No_{>0}
\]
if and only if \(\Lambda\) is a real multiple of constant-term extraction
\(\ct=[\omega^0]\).
\end{theorem}

\begin{proof}
Let \(w(\gamma)=\Lambda(\omega^\gamma)\).  Invariance gives
\(w(a\gamma)=w(\gamma)\) for all \(a>0\).  Any two positive exponents are
related by multiplication by a positive surreal, so \(w\) is constant on
\(\No_{>0}\); similarly it is constant on \(\No_{<0}\).  If the positive
constant were nonzero, the active set would contain the decreasing sequence
\(1>1/2>1/3>\cdots\), contradicting
\cref{prop:summation-finite-seq}.  If the negative constant were nonzero, it
would contain \(-1>-2>-3>\cdots\), again a contradiction.  Thus only the
weight at exponent \(0\) may survive, giving a multiple of \(\ct\).  The
converse is immediate because every \(\sigma_a\) fixes exponent \(0\) and its
coefficient.
\end{proof}

\begin{example}[Coefficient characters]
If \(d:(\No,+)\to(\R,+)\) is an additive class homomorphism, then
\(\chi_d(\gamma)=\exp(d(\gamma))\) is a positive character.  The resulting
automorphism fixes every monomial line and rescales it by a positive real.
A concrete choice is \(d=\ct\), constant-term extraction applied to the
\emph{exponent} \(\gamma\):
\[
  \omega^\gamma\longmapsto
  e^{\ct(\gamma)}\omega^\gamma.
\]
\end{example}

\subsection{Diagonal strong derivations}\label{subsec:derivations}

The same additive characters produce derivations rather than automorphisms.

\begin{theorem}[Diagonal derivations]\label{thm:diagonal-derivation}
Let \(d:(\No,+)\to(\R,+)\) be additive.  Then
\begin{equation}\label{eq:Dd}
 D_d\!\left(\sH_\gamma r_\gamma\omega^\gamma\right)
 =\sH_\gamma r_\gamma d(\gamma)\omega^\gamma
\end{equation}
is a strong \(\R\)-linear derivation of \(\No\):
\[
  D_d(xy)=xD_d(y)+yD_d(x).
\]
Conversely, every strong \(\R\)-linear derivation satisfying
\(D(\omega^\gamma)\in\R\omega^\gamma\) for every \(\gamma\) is uniquely of
this form.
\end{theorem}

\begin{proof}
The support of \(D_dx\) is contained in \(\supp(x)\), so the map is strong.
For monomials,
\begin{align*}
 D_d(\omega^{\gamma+\delta})
 &=d(\gamma+\delta)\omega^{\gamma+\delta}\\
 &=(d(\gamma)+d(\delta))\omega^{\gamma+\delta}\\
 &=\omega^\gamma D_d(\omega^\delta)
   +\omega^\delta D_d(\omega^\gamma).
\end{align*}
Hahn distributivity and local finiteness extend the Leibniz identity to all
normal forms.

Conversely, write
\(D(\omega^\gamma)=d(\gamma)\omega^\gamma\).  The Leibniz identity on a
product of two monomials forces
\(d(\gamma+\delta)=d(\gamma)+d(\delta)\).  Strong linearity determines \(D\)
on every normal form, giving \eqref{eq:Dd}.
\end{proof}

\begin{example}[A surreal Euler-type derivation]
With \(d=\ct\),
\[
  D_{\ct}(\omega^r)=r\omega^r\qquad(r\in\R),
\]
so in particular \(D_{\ct}(\omega)=\omega\), while
\(D_{\ct}\) kills real constants.  It is diagonal in the Conway--Hahn basis.
It is not asserted here to commute with the surreal exponential; differential
compatibility is an additional structure studied by Berarducci--Mantova and
later work \cite{BerarducciMantova,ADH}.
\end{example}

\subsection{Algebraic independence of monomials}

Linear coordinates also give a clean field-theoretic criterion.

\begin{theorem}[Rational exponents versus algebraic monomials]
\label{thm:monomial-alg-ind}
For \(\gamma_1,\ldots,\gamma_n\in\No\), the monomials
\(\omega^{\gamma_1},\ldots,\omega^{\gamma_n}\) are algebraically independent
over \(\R\) if and only if the exponents
\(\gamma_1,\ldots,\gamma_n\) are linearly independent over \(\Q\).
\end{theorem}

\begin{proof}
Suppose first that the exponents are \(\Q\)-dependent.  Clearing
denominators gives a nonzero integer relation
\(\sum_i m_i\gamma_i=0\).  Moving negative terms to the other side gives an
equality between two distinct monomials in the
\(\omega^{\gamma_i}\), hence a nonzero polynomial relation over \(\R\).

Conversely, assume the exponents are \(\Q\)-independent and let
\[
 P(X_1,\ldots,X_n)=\sum_{\mathbf{k}\in F}
 c_{\mathbf{k}}X_1^{k_1}\cdots X_n^{k_n}
\]
be a nonzero polynomial with finite exponent set
\(F\subseteq\N^n\).  Distinct multiindices \(\mathbf{k}\ne\mathbf{l}\)
give distinct surreal exponents
\[
  \sum_i k_i\gamma_i\ne\sum_i l_i\gamma_i,
\]
by rational independence.  Therefore
\[
 P(\omega^{\gamma_1},\ldots,\omega^{\gamma_n})
 =\sum_{\mathbf{k}\in F}
 c_{\mathbf{k}}\omega^{\sum_i k_i\gamma_i}
\]
is a nonzero finite normal form.  No cancellation between distinct exponents
is possible.
\end{proof}

\begin{corollary}
Every nontrivial monomial \(\omega^\gamma\), with \(\gamma\ne0\), is
transcendental over \(\R\).
\end{corollary}

\begin{corollary}[An explicit proper-class independent family]
\label{cor:proper-class-alg-independent}
The family
\begin{equation}\label{eq:proper-class-AI}
  \{\omega^{\omega^\alpha}:\alpha\in\On\}
\end{equation}
is algebraically independent over \(\R\).
\end{corollary}

\begin{proof}
The exponent family \(\{\omega^\alpha:\alpha\in\On\}\) is real-linearly,
hence rationally, independent: every finite relation would be a finite
normal form with distinct exponents \(\alpha\).  Apply
\cref{thm:monomial-alg-ind}.
\end{proof}

Thus the field extension \(\No/\R\) has proper-class transcendence complexity
in the evident class-theoretic sense.  An actual transcendence basis again
requires class choice and is not supplied canonically by normal form; the
explicit family \eqref{eq:proper-class-AI} gives a natural lower bound without
such a choice.

\section{Relation to the ProveIt collection and the literature}
\label{sec:repo-relation}

\subsection{Repository map}

The repository snapshot \texttt{\RepoShort} contains a broad surreal-number
program \cite{ProveIt}.  The present article is designed to complement, rather than repeat,
the following components.

\begin{center}
\small
\begin{tabularx}{0.98\textwidth}{L{0.34\textwidth}Y Y}
\toprule
Repository item & Principal object & Relation to this article\\
\midrule
\path{docs/surreal/vector-and-tensor-fields} & Finite-dimensional
\(K^n\), tensor fields, and support-controlled calculus over set-sized
surreal subfields & Uses surreal numbers as \emph{scalars}; does not ask for
a basis of \(\No\) over \(\R\).\\
\path{docs/surcomplex/three-duals-of-hahn-vector-spaces} &
\(V((t^\Gamma))\), algebraic scalar extension, valuation completion, strong
and continuous duals & Supplies the closest set-sized analogy to the strong
map and dual distinctions in \cref{sec:strong-maps}.\\
\path{docs/surcomplex/infinite-dimensional-hahn-spectral-theory} &
Hahn-summable vector families and Hahn-orthonormal synthesis & Explicitly
distinguishes Hahn bases from Hamel and Hilbert bases; the present report
applies that distinction to \(\No\) itself.\\
\path{docs/surcomplex/hahn-hilbert-geometry} & Set-sized Hahn--Hilbert
workspaces at surreal scales & Illustrates why analytical work selects a
set-sized workspace rather than the full fine class topology.\\
\path{docs/surreal/exponential-automorphism-rigidity} & Ordered,
valuation, and exponential automorphisms of \(\No\) & Complements the
monomial-preserving strong automorphisms in \cref{sec:multiplicative}; extra
exponential structure is substantially more rigid.\\
\bottomrule
\end{tabularx}
\end{center}

The repository states that its reports are AI-assisted, unrefereed drafts and
separately records which results have Lean formalizations.  The same claim
firewall is adopted here: no theorem newly proved in this report is presented
as machine checked merely because related foundations exist in the repository.

\subsection{Classical sources}

Conway introduced \(\No\) and its normal forms; Gonshor developed the sign-
sequence and omega-map theory in detail \cite{Conway,Gonshor}.  Hahn's
generalized series and Neumann's support lemma provide the ambient algebraic
technology \cite{Hahn,Neumann}.  Van den Dries and Ehrlich related the surreal
field and exponentiation to model-theoretic and Hahn-field structure
\cite{vdDE}; Ehrlich's surveys explain the role of \(\No\) as a universal
ordered continuum \cite{Ehrlich2012}.  Berarducci--Mantova and
Aschenbrenner--van den Dries--van der Hoeven developed the differential and
transseries aspects \cite{BerarducciMantova,ADH}.  Mantova--Matusinski give a
modern survey of Conway normal form and surreal analysis
\cite{MantovaMatusinski}.

The automorphism viewpoint used in \cref{sec:multiplicative} is now developed
systematically by Kaplan, Krapp, and Serra, who work in NBG with global choice,
view \(\No\) as a Hahn field, decompose its automorphism group, and isolate
strongly linear automorphisms \cite{KKS}.  Their work also notes that
choice-dependent bases enter descriptions of additive homomorphisms; this
reinforces the separation between intrinsic Hahn coordinates and arbitrary
Hamel coordinates.

The set-sized dimension formula in \cref{thm:window-dimension} is the
Erd\H{o}s--Kaplansky theorem.  We cite Jacobson's classical account and Bergman's
later discussion and extension \cite{Jacobson,Bergman}.

\subsection{What is genuinely specific to the vector-space question}

The literature generally invokes Conway normal form to study order, fields,
valuations, derivations, or automorphisms.  The present organization instead
starts from the linear question and follows its consequences:
\begin{enumerate}
\item it separates Hamel, Hahn, graded, and Schauder meanings of basis;
\item it proves the no-set-spanning theorem directly from support;
\item it quantifies the local Hamel-dimension explosion inside support
windows;
\item it exposes the quotient \(\No/\Fin\) by an explicit proper-class
independent family;
\item it derives an intrinsic strong real dual and contrasts it with the zero
order dual;
\item it explains why the full proper-class valuation topology cannot realize
normal forms as set-indexed limits.
\end{enumerate}
These are primarily structural deductions from classical coordinates.  Their
value is to prevent category errors when importing finite-dimensional or
Banach-space intuitions into surreal analysis.

\section{Further research questions and proposed program}
\label{sec:research}

The vector-class viewpoint opens a research program that is distinct from,
but compatible with, the existing surreal analysis, automorphism, and
formalization programs.  The questions below are ordered roughly from
foundational to analytic and computational.

\subsection{Foundations and definability}

\begin{question}[Choice strength for class bases]
What is the weakest conservative class theory over ZFC that proves every
class vector space whose carrier is bijective with \(\On\) has a Hamel class
basis?  In particular, how much of global choice is needed for
\cref{thm:class-basis}, and can the statement fail in natural models of NBG
without global choice?
\end{question}

The staged proof uses global choice to select basis extensions through a
proper class of stages.  A reverse-mathematical analysis would separate the
choice needed for this particular highly structured vector class from a
general class-basis principle.

\begin{question}[Definable Hamel bases]
Does \(\No\) admit a Hamel basis definable, without arbitrary parameters, in
natural class universes such as \((V,\in)\), \(L\), or HOD?  Can one prove
that no such basis is invariant under specified large subgroups of
\(\Aut(\No,+,\cdot,<)\)?
\end{question}

A useful intermediate target is a definable basis containing all monomials,
as in \cref{prop:basis-containing-monomials}, together with a definable rule
for selecting the necessary infinite-support supplement.  Symmetry may make
this impossible in the full universe even when a basis exists abstractly.

\begin{question}[Class matroids]
Develop a class-matroid framework in which linear independence on \(\No\),
basis extension, closure, and rank are formulated without forcing all ranks
to be set cardinals.  Which ordinary matroid theorems survive for set-finitary
class matroids under global choice?
\end{question}

The linear matroid of \(\No\) is finitary---dependence is witnessed by a
finite subset---but its ground class and every basis are proper.  This makes
it a natural test case for class-sized combinatorial geometry.

\subsection{Support strata and quotient geometry}

\begin{question}[Structure of \(\No/\Fin\)]
Classify the real vector-class quotient \(\No/\Fin\) and its natural
filtrations by support order type, cofinality, and leading exponent.  Is there
a canonical hierarchy of subquotients whose set-sized pieces have computable
Hamel dimensions?
\end{question}

The block family in \cref{thm:quotient-independent} gives only a lower bound.
A finer analysis could distinguish countable support, bounded support order
type, and support contained in specified exponent intervals.

\begin{question}[Support-rank dimension spectra]
For an infinite cardinal \(\kappa\), let \(\No^{(<\kappa)}\) consist of
surreals whose normal-form support has cardinality below \(\kappa\), or order
type below \(\kappa\).  Determine the Hamel dimensions of set-sized
intersections of these classes and the strictness of the resulting hierarchy.
\end{question}

The countable-window theorem suggests cardinal-exponent formulas, but order
type and ambient exponent geometry can interact when the allowed support
window itself is not reverse well ordered.

\begin{question}[Canonical complements]
Which naturally defined real subspaces of \(\No\) admit canonical strong
complements?  Support-defined subspaces do by \cref{thm:support-splitting};
what happens for birthday cutoffs, differential kernels, logarithmic strata,
or fields generated by selected monomials?
\end{question}

The contrast between support subspaces and birthday subspaces may reveal
which decompositions belong to the Hahn geometry and which depend on the
genetic construction.

\subsection{Strong operator algebras and duality}

\begin{question}[The algebra of Hahn-compatible matrices]
Give intrinsic necessary and sufficient conditions for composition,
invertibility, triangular factorization, and spectral decomposition of the
class matrices in \cref{thm:strong-matrix}.  Identify natural operator ideals:
finite-row, finite-column, support-contracting, valuation-increasing, and
nuclear analogues.
\end{question}

A promising route is to first solve the problem on every set-sized support
window and then formulate uniform class conditions.  The recent theory of
automorphisms and derivations on algebras with formal infinite sums should be
compared directly \cite{BagayokoEtAl}.

\begin{question}[Invariant strong duals]
For important automorphism groups \(G\) of \(\No\), compute
\[
  (\No^{\vee}_{\mathrm{str}})^G.
\]
Theorem~\ref{thm:dilation-invariant-dual} shows that invariance under all
positive exponent dilations leaves exactly the constant-term line.  What are
the invariant strong functionals for valuation-preserving, omega-preserving,
or exponential automorphism groups?
\end{question}

This is a tractable bridge between the linear theory developed here and the
automorphism decomposition of Kaplan--Krapp--Serra.

\begin{question}[Duality beyond real scalars]
Classify strong, order-bounded, and topology-continuous maps from \(\No\) to
other codomains: a set-sized Hahn field, \(\No\) itself, the surcomplex field,
or quotient spaces such as \(\No/\Fin\).  Which codomains avoid the collapse
of \cref{thm:order-dual-zero}?
\end{question}

Real-valued positive states vanish because the codomain is Archimedean.
Surreal-valued positive maps may exist and could provide a better foundation
for surreal measure, probability, and convex analysis.

\subsection{Topology and generalized bases}

\begin{question}[Coarser topologies with genuine Schauder expansions]
Classify set-generated valuation or bornological structures on \(\No\) for
which selected classes of Conway normal forms are limits of their finite
truncations.  What is the largest natural subspace on which the monomials form
a Schauder basis for a given set-sized scale system?
\end{question}

Any answer must evade \cref{thm:eventually-constant} by replacing the full
proper-class neighborhood basis with a set-sized cofinal family, restricting
the domain, or allowing class-indexed convergence.

\begin{question}[Class-indexed convergence]
Can class-indexed nets or filters be developed safely enough to make every
Conway normal form a topological sum without collapsing separation axioms or
leaving conservative class theory?  How does such convergence compare with
strong Hahn summation?
\end{question}

Strong summation may already be the correct replacement.  A topological
reconstruction would be valuable only if it adds robust continuity,
completion, or functional-analytic theorems.

\begin{question}[Bornological and locally convex analogues]
Is there a useful class-bornology on \(\No\) whose bounded linear dual lies
strictly between the strong dual and the zero order-bounded real dual?  Can
one formulate barrelledness, nuclearity, or reflexivity using only set-sized
bounded pieces?
\end{question}

The support windows \(\Hahn(A)\) provide a canonical directed family of
set-sized test spaces.  An inductive-limit or bornological approach may be
more faithful than the full fine topology.

\subsection{Differential and exponential compatibility}

\begin{question}[Derivation-adapted bases]
How does the Berarducci--Mantova derivation act relative to the monomial Hahn
basis, and is there a more refined basis---perhaps using log-atomic numbers or
transseries trees---in which its matrix is triangular or locally finite?
\end{question}

The diagonal derivations of \cref{thm:diagonal-derivation} are completely
controlled by additive real characters, but the canonical surreal derivation
changes scales in a subtler way.  A basis adapted simultaneously to
multiplication, exponentiation, and derivation would be highly significant.

\begin{question}[Exponential rigidity of monomial automorphisms]
Determine exactly which pairs \((\chi,\phi)\) in
\cref{thm:monomial-aut} commute with the surreal exponential or preserve the
omega map.  Existing rigidity results rule out broad subclasses; an explicit
intersection with the monomial automorphism group should be computable.
\end{question}

A direct coefficient-and-leading-scale proof could provide an independent
check on more general automorphism theorems.

\begin{question}[Transcendence bases compatible with support]
Can one construct a class transcendence basis of \(\No/\R\) whose elements
have controlled normal-form support, and then describe the algebraic closure
of the generated field inside \(\No\)?  How far is the explicit family
\(\{\omega^{\omega^\alpha}\}\) from a transcendence basis?
\end{question}

The answer should distinguish finite field generation from Hahn closure;
normal-form completeness is far larger than rational-function generation.

\subsection{Tensor constructions and geometry}

\begin{question}[Completed tensor powers]
Define and compare the algebraic tensor powers of the class vector space
\(\No\) with support-completed Hahn tensor products.  Do the monomials
\(\omega^{\gamma_1}\otimes\cdots\otimes\omega^{\gamma_n}\) form a natural
multi-Hahn basis, and which support conditions are forced by contraction and
symmetrization?
\end{question}

This would connect the present scalar decomposition with the repository's
finite-dimensional tensor-field program and with infinite-dimensional Hahn
Hilbert geometry.

\begin{question}[Exterior and symmetric class algebras]
Describe the exterior, symmetric, and Clifford algebras of \(\No\) over
\(\R\) under both Hamel and Hahn completions.  Which constructions retain a
canonical grading by surreal exponents?
\end{question}

The algebraic versions depend on a Hamel basis only through isomorphism, while
support-completed versions may preserve much more of the normal-form geometry.

\subsection{Formalization and computation}

\begin{question}[Formalization strategy]
Formalize the set-sized core of the article in Lean: reverse-well-ordered
support spaces, strong sums, support projections, the matrix theorem, the
strong-dual criterion, the Vandermonde lower bound, and the monomial algebraic-
independence theorem.  Then isolate the additional class-theoretic axioms
needed for global statements.
\end{question}

Most proofs reduce to finite-support bookkeeping plus well-order properties
and are good candidates for machine checking.  The class-basis existence
theorem should be kept logically separate from constructive normal-form
results.

\begin{question}[Symbolic linear algebra for normal forms]
Develop algorithms that operate on finite or certified transfinite
truncations: support projection, leading-grade elimination, diagonal and
reindexing operators, strong matrix application, and independence tests for
monomial families.  Produce proof certificates for all support conditions.
\end{question}

Such tools would turn the Hahn-basis viewpoint into a practical computer
algebra interface.  A certificate should state the support order type, local
finiteness bounds, and the exact truncation region used.

\begin{question}[Repository integration]
Add a dedicated ProveIt report and Lean ledger entry that distinguishes
\(\No^n\) over \(\No\), set-sized Hahn vector spaces, and \(\No\) over
\(\R\).  Cross-link the terminology ``Hamel basis,'' ``Hahn basis,''
``strongly summable family,'' and ``associated graded'' across all related
reports.
\end{question}

This editorial step would prevent the same word ``basis'' from denoting three
incompatible synthesis rules in neighboring parts of the project.

\section{Conclusion}\label{sec:conclusion}

The surreal numbers are unquestionably a vector class over \(\R\), but the
ordinary phrase ``choose a basis'' suppresses the structure that makes them
surreal.  A Hamel basis sees only finite sums.  It exists under global choice,
may be chosen to contain every monomial, and is necessarily supplemented by a
proper class of infinite-support vectors.  Such a basis is useful for abstract
algebraic existence arguments but not for normal-form computation, order, or
asymptotics.

The natural answer is instead
\[
  \boxed{\M=\{\omega^\gamma:\gamma\in\No\}}.
\]
This class is canonical in three mutually reinforcing senses:
\begin{enumerate}
\item it is the Conway--Hahn basis giving every surreal a unique strongly
summable real expansion;
\item it is the basis of the one-dimensional dominance grades, with
\(\gr(\No)\cong\R[\No_{\mathrm{add}}]\);
\item it is normalized by simplicity, since \(\omega^\gamma\) is the
simplest positive representative of its Archimedean class.
\end{enumerate}
Its basis law is therefore not
\(\No=\bigoplus_\gamma\R\omega^\gamma\), but the sharper statement
\[
  \No=\classprod_{\gamma\in\No}\R\omega^\gamma,
\]
where the product sign denotes set-supported, reverse-well-ordered Hahn
coordinates.  (Equivalently, use \eqref{eq:direct-vs-hahn}; the displayed
formula is conceptual rather than new notation.)

Several consequences follow immediately once the correct synthesis rule is
fixed.  Support subsets give canonical projections.  Strong operators have
Hahn matrices.  Strong real functionals are precisely summation-finite
weighted coefficient sums, while order-bounded real functionals vanish.
Set-sized support windows reveal an enormous gap between Hahn and Hamel
dimensions.  The full fine topology is too refined for set-indexed partial
sums, so normal forms are formal strong sums rather than Schauder series.
Monomial symmetries and diagonal derivations reduce to group maps on the
exponents, and rational exponent independence becomes algebraic monomial
independence.

The resulting picture is analogous to several familiar distinctions, but is
not identical to any one of them: Fourier series versus finite trigonometric
polynomials, formal power series versus polynomials, completed graded objects
versus their direct sums, and topological bases versus algebraic bases.  The
surreal case adds two features absent from those set-sized analogies: the
exponent group is itself the proper class \(\No\), and every individual
series still has only set support.  That combination is exactly why the
Conway--Hahn basis is both global and rigorously manageable.

\appendix

\section{Proof and dependency audit}\label{app:audit}

\begin{center}
\small
\begin{longtable}{L{0.26\textwidth}L{0.26\textwidth}L{0.36\textwidth}}
\toprule
Result & Main input & Status in this report\\
\midrule
\endfirsthead
\toprule
Result & Main input & Status in this report\\
\midrule
\endhead
Conway--Hahn basis, \cref{thm:hahn-basis} & Conway normal-form theorem &
Classical theorem restated in vector-language.\\
No set spans, \cref{thm:no-set-spans} & Set union of normal-form supports &
Complete direct proof.\\
Hamel class basis, \cref{thm:class-basis} & NBG+GC; set-sized basis extension &
Complete class-recursion proof.\\
Basis containing monomials,
\cref{prop:basis-containing-monomials} & NBG+GC; greedy class recursion &
Complete proof.\\
Quotient independent family,
\cref{thm:quotient-independent} & Disjoint countable support blocks &
Complete explicit proof.\\
Support-window dimension,
\cref{thm:window-dimension} & Erd\H{o}s--Kaplansky theorem & Classical cardinal
result applied to \(\R^A\); countable lower bound also proved directly.\\
Support projections, \cref{thm:support-splitting} & Heredity of reverse well
ordering & Complete proof.\\
Associated graded, \cref{thm:graded} & Leading-term multiplication & Complete
proof.\\
Strong matrix classification,
\cref{thm:strong-matrix} & Flattening of locally finite Hahn families &
Complete proof.\\
Strong real dual, \cref{thm:strong-dual} & Strong matrix theorem with
codomain \(\R\) & Complete proof.\\
Dilation invariants,
\cref{thm:dilation-invariant-dual} & Strong-dual weights and exponent
dilations & Complete proof.\\
Order dual is zero,
\cref{thm:order-dual-zero} & Infinite surreal bound \(\omega\abs{x}\) &
Complete proof.\\
Set-indexed nets eventually constant,
\cref{thm:eventually-constant} & Every set of surreals has a surreal upper
bound & Complete proof.\\
Monomial automorphisms,
\cref{thm:monomial-aut} & Hahn support transport and character laws & Complete
proof for the stated monomial-line-preserving subgroup.\\
Diagonal derivations,
\cref{thm:diagonal-derivation} & Additive real characters of the exponent
group & Complete proof.\\
Monomial algebraic independence,
\cref{thm:monomial-alg-ind} & Uniqueness of finite normal forms & Complete
proof.\\
\bottomrule
\end{longtable}
\end{center}

No computer algebra was needed for these proofs.  The main verification risks
are foundational rather than computational: distinguishing sets from proper
classes, and never replacing strong summability by an unjustified topological
limit.

\section{Notation and terminology}\label{app:notation}

\begin{center}
\small
\begin{tabularx}{0.96\textwidth}{L{0.24\textwidth}Y}
\toprule
Symbol & Meaning\\
\midrule
\(\No\), \(\On\) & The surreal-number and ordinal classes.\\
\(\omega^\gamma\) & Conway omega-map monomial at exponent \(\gamma\).\\
\(\supp(x)\) & Set of exponents with nonzero coefficient in the Conway normal
form of \(x\).\\
\([\omega^\gamma]x\) & Real coefficient of \(\omega^\gamma\) in \(x\).\\
\(\sH\) & Strong or Hahn sum.\\
\(\M\) & The monomial class \(\{\omega^\gamma:\gamma\in\No\}\).\\
\(\Fin\) & Finite-support normal forms; the algebraic direct sum of the
monomial lines.\\
\(\Hahn(A)\) & Surreals whose support is contained in the reverse-well-ordered
set \(A\).\\
\(P_A\) & Projection retaining exactly the exponents in \(A\).\\
\(v(x)\) & Natural valuation \(-\lead(x)\).\\
Strong map & A real-linear map preserving every admissible set-indexed Hahn
sum and commuting with it.\\
Summation-finite class & A class meeting every reverse-well-ordered set in a
finite set; equivalently, containing no infinite decreasing sequence.\\
\bottomrule
\end{tabularx}
\end{center}

\clearpage
\begin{thebibliography}{99}

\bibitem{ADH}
M. Aschenbrenner, L. van den Dries, and J. van der Hoeven,
\emph{Asymptotic Differential Algebra and Model Theory of Transseries},
Annals of Mathematics Studies 195, Princeton University Press, 2017.

\bibitem{BagayokoEtAl}
V. Bagayoko, L. S. Krapp, S. Kuhlmann, D. Panazzolo, and M. Serra,
``Automorphisms and derivations on algebras endowed with formal infinite
sums,'' arXiv:2403.05827, revised 2025.

\bibitem{BerarducciMantova}
A. Berarducci and V. Mantova,
``Surreal numbers, derivations and transseries,''
\emph{Journal of the European Mathematical Society} 20 (2018), 339--390.
\href{https://doi.org/10.4171/JEMS/769}{doi:10.4171/JEMS/769}.

\bibitem{Bergman}
G. M. Bergman,
``Families of ultrafilters, and homomorphisms on infinite direct product
algebras,'' arXiv:1301.6383, 2013; see Section 6 for the
Erd\H{o}s--Kaplansky theorem and reduced products.

\bibitem{Conway}
J. H. Conway,
\emph{On Numbers and Games}, 2nd ed., A K Peters, 2001.

\bibitem{Ehrlich2012}
P. Ehrlich,
``The absolute arithmetic continuum and the unification of all numbers great
and small,'' \emph{Bulletin of Symbolic Logic} 18 (2012), no. 1, 1--45.
\href{https://doi.org/10.2178/bsl/1327328438}{doi:10.2178/bsl/1327328438}.

\bibitem{Gonshor}
H. Gonshor,
\emph{An Introduction to the Theory of Surreal Numbers},
London Mathematical Society Lecture Note Series 110,
Cambridge University Press, 1986.
\href{https://doi.org/10.1017/CBO9780511629143}{doi:10.1017/CBO9780511629143}.

\bibitem{Hahn}
H. Hahn,
``{\"U}ber die nichtarchimedischen Gr{\"o}{\ss}ensysteme,''
\emph{Sitzungsberichte der Kaiserlichen Akademie der Wissenschaften,
Mathematisch-Naturwissenschaftliche Klasse} 116 (1907), 601--655.

\bibitem{Jacobson}
N. Jacobson,
\emph{Lectures in Abstract Algebra, Volume II: Linear Algebra},
Van Nostrand, 1953; reprinted as Graduate Texts in Mathematics 31,
Springer, 1975.

\bibitem{KKS}
E. Kaplan, L. S. Krapp, and M. Serra,
``Decomposing the automorphism group of the surreal numbers,''
arXiv:2509.22374v3, revised April 23, 2026.
\href{https://doi.org/10.48550/arXiv.2509.22374}{doi:10.48550/arXiv.2509.22374}.

\bibitem{MantovaMatusinski}
V. Mantova and M. Matusinski,
``Surreal numbers with derivation, Hardy fields and transseries: a survey,''
arXiv:1608.03413v2, 2016.

\bibitem{MO2013}
M. Battaglia, with comments by T. Trimble,
``Is a certain group, derivable from the surreal numbers, isomorphic to the
surreal numbers?,'' MathOverflow question 140732, August--September 2013.
\url{https://mathoverflow.net/questions/140732/}.

\bibitem{Neumann}
B. H. Neumann,
``On ordered division rings,''
\emph{Transactions of the American Mathematical Society} 66 (1949),
202--252.

\bibitem{ProveIt}
V. Reshetnikov,
\emph{ProveIt}, GitHub repository, surreal-number documentation and Lean
sources, snapshot \texttt{\RepoCommit}, accessed October 1, 2026.
\url{https://github.com/VladimirReshetnikov/ProveIt}.

\bibitem{vdDE}
L. van den Dries and P. Ehrlich,
``Fields of surreal numbers and exponentiation,''
\emph{Fundamenta Mathematicae} 167 (2001), no. 2, 173--188;
erratum, 168 (2001), 295--297.

\end{thebibliography}

\end{document}
